\scanpage{208}
\mspage{82}
Ceci posé, on voit donc, $C$ étant $\unsure{\pi|R}$-accessible et $I$ grand devant $\pi$, que l'ens. $I_C$ des $i\in I$ tels que $f_i\in C$ est \uline{cofinal} dans $I$. Mais \deleted{les $f_i$} comme $f_i\in\mathrm{Filt}^\pi$, $f_i\in C$ \unsure{revient|équivaut} à $f_i\in C^\pi$. Comme $A\unsure{-|\longrightarrow}B$ n'est pas iso, il existe un $i_0$ tel que $B_{i_0}\not\subset A$, \deleted{\unsure{donc}} \inserted{puis} un $i\geq i_0$ tel que $f_i\in C^\pi$, \inserted{et aussi} $B_i\not\subset A$ (puisque $B_i\supset B_{i_0}$). Ainsi on a prouvé ceci

$(\ast)$ Sous les hyp. de (6.4.) et moyennant l'hyp. \textcircled{e'} sur $M$, il existe un sous-objet $B'$ de $B$ satisfaisant les conditions suivantes

a) $B'\in\mathrm{Filt}^\pi$

b) $B'\not\subset A$

c) Posant $A'=B'\cap A$, l'inclusion
\[
 i':A'-B'
\]
est dans $C^\pi$.

De plus, l'ens des $J\subset\mathrm{Ss}_\pi(B)$ formé des $B'$ en question est cofinal dans $\mathrm{Ss}_\pi(B)$, et $=$ $\pi$-cofinal.

Il faut prouver encore ce dernier point --- que $J$ (\deleted{\unsure{dont}} on sait qu'il est cofinal) est $=$ $\pi$-cofinal, \deleted{\unsure{il est donc}} i.e. stable par $\sup$ \inserted{(dans $\mathrm{Ss}_\pi(B)$)} de \deleted{\unsure{ses}} parties filtrantes $J'\in J$ telles que $\operatorname{card}J'\leq\pi$. Mais cela résulte aussitôt du fait que $C$ étant stable par
% archive sheet number: 207
